\scanpage{13}
Soit $W$ une variété \deleted{topologique} diff.

$P,Q$ deux sous-variétés à bord

avec
\[
P\cap Q=\partial P=\partial Q
\]
\[
P\cup Q=W
\]
et $P$ et $Q$ étant \unsure{deux|des} fibrés diff.

\deleted{Alors on a}
\margin{$\partial P$}
\margin{\deleted{\unsure{W}}}
\diagram{\begin{tikzcd}[column sep=small]
P\arrow[r]&S&Q\arrow[r]&T\\
M\arrow[u,hook]&&N\arrow[u,hook]
\end{tikzcd}}
de bases $S,T$ des variétés diff \unsure{cptes} $S,T$, de fibres des variétés à bords $M$, resp. $N$. Donc $\partial P$ est un fibré sur $S$ de fibre $\partial M$, et $\partial Q$ un fibré sur $T$ de fibre $\partial N$ :
\diagram{\begin{tikzcd}[column sep=small]
\partial P\arrow[r]&S&\partial Q\arrow[r]&T\\
\partial M\arrow[u,hook]&&\partial N\arrow[u,hook]
\end{tikzcd}}
On suppose donnés des isomorphismes
\[
\alpha:S\simeq\partial N\qquad\beta:T\simeq\partial M
\]
et \deleted{pour les fibres} des trivialisations
\[
\sigma:\partial P\simeq S\times\partial M\qquad\tau:\partial Q\simeq T\times\partial N
\]
D'où, par composition
\[
(\alpha,\sigma):\partial P\simeq\partial N\times\partial M,\quad(\beta,\tau):\partial Q\simeq\partial M\times\partial N
\]
et par suite un isom
\[
\varphi:\partial P\simeq\partial Q
\]
\mspage{12}
